\UseRawInputEncoding

\documentclass[
% -- opções da classe memoir --
article,			% indica que é um artigo acadêmico
12pt,				% tamanho da fonte
oneside,			% para impressão apenas no verso. Oposto a twoside
a4paper,			% tamanho do papel.
hidelinks, 
% -- opções da classe abntex2 --
%chapter=TITLE,		% títulos de capítulos convertidos em letras maiúsculas
%section=TITLE,		% títulos de seções convertidos em letras maiúsculas
%subsection=TITLE,	% títulos de subseções convertidos em letras maiúsculas
%subsubsection=TITLE % títulos de subsubseções convertidos em letras maiúsculas
% -- opções do pacote babel --
english,			% idioma adicional para hifenização
brazil,				% o último idioma é o principal do documento
]{abntex2}

% ---
% Pacotes fundamentais 
% ---
\usepackage{lmodern}			% Usa a fonte Latin Modern
\usepackage[T1]{fontenc}		% Selecao de codigos de fonte.
\usepackage[utf8]{inputenc}		% Codificacao do documento (conversão automática dos acentos)
\usepackage{indentfirst}		% Indenta o primeiro parágrafo de cada seção.
\usepackage{nomencl} 			% Lista de simbolos
\usepackage{color}				% Controle das cores
\usepackage{graphicx}			% Inclusão de gráficos
\usepackage{microtype} 			% para melhorias de justificação
\usepackage{amsmath}			% alinhar as equações
\usepackage{amssymb}
\usepackage{amsthm}
\usepackage[brazilian,noabbrev]{cleveref}
\usepackage{multirow}			
\usepackage{tikz}				% Grafico
\usetikzlibrary{calc}			% alinhar o grafico por coordenadas
\usepackage{physics}
\usepackage{caption}
\usepackage{tikz-cd}
\usepackage{threeparttable}
\usepackage{threeparttablex}
% ---


% Ambiente de teoremas e provas
\renewenvironment{proof}{{\bfseries \textit{Prova}.	}}{\qed}

\theoremstyle{plain}
\newtheorem{teorema}{Teorema}

\theoremstyle{definition}
\newtheorem{definicao}{Definição}

\theoremstyle{remark}
\newtheorem{rascunho}{Rascunho}[teorema]
% ---


% Comandos criados
\newcommand{\R}{\mathbb{R}}   % comando para facilitar ao escrever o R de numeros reais - \R 
\newcommand\inner[2]{\left \langle #1, #2 \right \rangle}   % Comando para criar um porduto interno
\crefformat{equation}{equação #2#1#3}
\crefmultiformat{equation}{equações #2#1#3}{ e #2#1#3}{, #2#1#3}{ e #2#1#3}
% ---





% Pacotes tabela
\usepackage{booktabs,tabularx}
\usepackage[input-decimal-markers=,]{siunitx}
\usepackage{float}
\usepackage{longtable}
% ---

% ---
% Pacotes de citações
% ---
\usepackage[brazilian,hyperpageref]{backref}	 % Paginas com as citações na bibl
\usepackage[alf, abnt-emphasize=bf,abnt-etal-cite=2,abnt-etal-list=2]{abntex2cite}	% Citações padrão ABNT
% ---



% ---
% Informações de dados para CAPA e FOLHA DE ROSTO
% ---
\instituicao{a}
\tituloestrangeiro{}
\titulo{Choques monetário nos preços dos ativos}
\autor{Gabriel Arruda\thanks{Doutorando em Economia - PPGEco-UFSC} \\ g.arruda@hotmail.com}
\data{}
% ---

% ---
% Configurações de aparência do PDF final

% alterando o aspecto da cor azul
\definecolor{blue}{RGB}{41,5,195}

% informações do PDF
\makeatletter
\hypersetup{
	%pagebackref=true,
	pdftitle={\@title}, 
	pdfauthor={\@author},
	pdfsubject={Macroeconomia},
	pdfcreator={Gabriel Arruda},
	pdfkeywords={SDFM}{choques monetários}{mercado financeiro}{macroeconomia}{politica monetária}, 
	colorlinks=true,       		% false: boxed links; true: colored links
	linkcolor=black,          	% color of internal links
	citecolor=black,        		% color of links to bibliography
	filecolor=magenta,      		% color of file links
	urlcolor=black,
	bookmarksdepth=4
}
\makeatother

% ---
% Altera as margens padrões
% ---
\setlrmarginsandblock{3cm}{3cm}{*}
\setulmarginsandblock{3cm}{3cm}{*}
\checkandfixthelayout
% ---

% O tamanho do parágrafo é dado por:
\setlength{\parindent}{1.3cm}

% Controle do espaçamento entre um parágrafo e outro:
\setlength{\parskip}{0.2cm}  % tente também \onelineskip

% Espaçamento simples
\SingleSpacing

% ----
% Início do documento
% ----
\begin{document}

% Retira espaço extra obsoleto entre as frases.
\frenchspacing 

% ----------------------------------------------------------
% ELEMENTOS PRÉ-TEXTUAIS
% ----------------------------------------------------------

%---
%
% Se desejar escrever o artigo em duas colunas, descomente a linha abaixo
% e a linha com o texto ``FIM DE ARTIGO EM DUAS COLUNAS''.
% \twocolumn[    		% INICIO DE ARTIGO EM DUAS COLUNAS
%
%---
% página de titulo
\maketitle

% resumo em português
\renewcommand{\resumoname}{Resumo}
\begin{resumoumacoluna}
    Um choque contracionista de 50 pontos-base na política monetária brasileira leva a uma queda imediata de aproximadamente 3\% no mercado acionário e a uma apreciação de 8\% na taxa de câmbio. A análise também revela um aumento significativo nos yields de títulos de diferentes maturidades e respostas assimétricas entre as classes de ativos, com efeitos mais pronunciados nos ativos de maior risco e no setor imobiliário. Estes resultados são obtidos por meio de um modelo de fatores dinâmicos estruturais (SDFM), que supera as limitações de não-fundamentalidade dos modelos VAR tradicionais ao incorporar um amplo conjunto de informações macroeconômicas e financeiras do Brasil.
\vspace{\onelineskip} 
\end{resumoumacoluna}


\textual

\section{Intro}

A compreensão dos canais de transmissão da política monetária é um dos temas centrais da macroeconomia, pois as decisões do banco central afetam diretamente as condições de crédito, as expectativas dos agentes e a atividade econômica. Dentre os diversos canais, a resposta dos preços dos ativos financeiros — como ações, títulos e a taxa de câmbio — é particularmente relevante, pois reflete a reação imediata dos mercados e serve como um importante mecanismo de transmissão para a economia real. A literatura recente tem dedicado atenção considerável a quantificar essa relação, utilizando modelos cada vez mais sofisticados para isolar o impacto de choques monetários na precificação de risco (e.g., \citeonline{BEKAERT2013771}, \citeonline{alessi}). No contexto de economias emergentes como o Brasil, caracterizadas por maior volatilidade e particularidades institucionais, a análise desses efeitos torna-se ainda mais crucial para a condução da política monetária e a gestão de portfólios.

A correta identificação desses efeitos depende crucialmente da informação disponível para o econometrista. Modelos tradicionais, como os Vetores Autorregressivos Estruturais (SVAR), frequentemente sofrem do problema da não-fundamentalidade, pois não conseguem incorporar o vasto conjunto de dados que os agentes econômicos e o próprio banco central utilizam em suas decisões, o que pode levar a uma má especificação do choque monetário \cite{villaverde}. Estudos mostram que, por consequência, os SVARs tendem a subestimar a magnitude do impacto da política monetária e a gerar respostas contraintuitivas em algumas variáveis \cite{alessi}.

Para superar essa limitação, este estudo emprega um Modelo de Fatores Dinâmicos Estruturais (SDFM). Ao extrair fatores comuns de um grande conjunto de variáveis macroeconômicas e financeiras, o SDFM permite uma identificação mais robusta dos choques, capturando a informação relevante de maneira parcimoniosa \cite{STOCK2016415}. A principal contribuição deste trabalho é, portanto, fornecer uma quantificação mais precisa dos canais de transmissão da política monetária no Brasil, mostrando que os efeitos sobre os preços dos ativos são mais fortes e rápidos do que a literatura baseada em modelos de menor dimensão tipicamente encontra. O restante do artigo detalha a metodologia, os dados e aprofunda a análise dos resultados.

\section{Revisão de literatura}

A literatura sobre os efeitos da política monetária nos mercados financeiros tem evoluído significativamente nas últimas décadas. Esta evolução é marcada tanto pelo desenvolvimento de modelos teóricos cada vez mais sofisticados quanto pelo surgimento de novas metodologias empíricas que permitem uma melhor identificação dos choques monetários e seus efeitos \cite{STOCK2016415}.

Os fundamentos teóricos desta literatura são estabelecidos por modelos de equilíbrio geral que buscam compreender os mecanismos de transmissão da política monetária para os mercados financeiros. \citeonline{castelnuovo} desenvolvem um modelo DSGE que incorpora rigidez nominal de salários e interações entre diferentes tipos de investidores, demonstrando como flutuações na riqueza financeira afetam o consumo. Complementarmente, \citeonline{Cooley} propõem um modelo de equilíbrio geral com empresas heterogêneas que revela como a estrutura financeira das empresas influencia sua resposta a choques monetários.

Embora os modelos teóricos forneçam importantes insights, a necessidade de quantificar precisamente estes efeitos levou ao desenvolvimento de análises empíricas. Os primeiros trabalhos utilizaram modelos VAR tradicionais. \citeonline{CHATZIANTONIOU2013754} empregam um VAR para mostrar como as políticas fiscal e monetária afetam conjuntamente o desempenho do mercado acionário. No entanto, os modelos VAR tradicionais enfrentam limitações na identificação precisa dos choques monetários. Esta limitação motivou o desenvolvimento dos modelos SVAR, que impõem restrições teoricamente fundamentadas para identificar os choques estruturais. \citeonline{BEKAERT2013771} decompõem o índice VIX em componentes de aversão ao risco e incerteza, e utilizam um SVAR para mostrar como a política monetária expansionista afeta a aversão ao risco.

Uma limitação importante dos modelos VAR/SVAR é sua incapacidade de incorporar grandes conjuntos de informação, o que pode levar ao problema de não-fundamentalidade dos choques estruturais. \citeonline{alessi} fornecem evidência clara deste problema ao comparar os resultados de um SVAR com um modelo de fatores dinâmicos estrutural (SDFM) usando o mesmo instrumento de alta frequência para identificar choques monetários. Os autores mostram que os SVARs subestimam consistentemente a magnitude do impacto da política monetária e produzem respostas contra-intuitivas em várias variáveis financeiras. Em contraste, o SDFM, ao incorporar um conjunto mais amplo de informações, encontra efeitos mais fortes e mais rápidos nos preços dos ativos, consistentes com a teoria macroeconômica padrão e com evidências de estudos de eventos. Outros modelos de grande dimensão, como os FAVAR, também têm contribuído para esta literatura \cite{lutz,miranda}. Metodologias mais recentes buscam refinar ainda mais a identificação, utilizando, por exemplo, instrumentos de alta frequência \cite{miranda-transmi,alessi} ou modelos com parâmetros variáveis no tempo \cite{pascal}.


% \section{Revisão de literatura}

% A literatura sobre os efeitos da política monetária nos mercados financeiros tem evoluído significativamente nas últimas décadas. Esta evolução é marcada tanto pelo desenvolvimento de modelos teóricos cada vez mais sofisticados quanto pelo surgimento de novas metodologias empíricas que permitem uma melhor identificação dos choques monetários e seus efeitos \cite{STOCK2016415}.

% Os fundamentos teóricos desta literatura são estabelecidos por modelos de equilíbrio geral que buscam compreender os mecanismos de transmissão da política monetária para os mercados financeiros. \citeonline{castelnuovo} desenvolvem um modelo DSGE que incorpora rigidez nominal de salários e interações entre diferentes tipos de investidores, demonstrando como flutuações na riqueza financeira afetam o consumo. O modelo fornece evidências robustas do papel dos preços das ações na atividade real e identifica uma resposta significativa do Federal Reserve às flutuações do mercado acionário. Complementarmente, \citeonline{Cooley} propõem um modelo de equilíbrio geral com empresas heterogêneas que revela como a estrutura financeira das empresas influencia sua resposta a choques monetários. Um resultado central deste modelo é que empresas menores, devido à maior alavancagem financeira, são mais sensíveis a variações nas taxas de juros.

% Embora os modelos teóricos forneçam importantes insights sobre os mecanismos de transmissão, a necessidade de quantificar precisamente estes efeitos levou ao desenvolvimento de análises empíricas. Os primeiros trabalhos empíricos utilizaram modelos VAR tradicionais. \citeonline{CHATZIANTONIOU2013754} empregam um VAR para mostrar como as políticas fiscal e monetária afetam conjuntamente o desempenho do mercado acionário. Seus resultados revelam diferenças importantes entre países: enquanto no Reino Unido tanto a política fiscal quanto a monetária afetam diretamente o mercado de ações, na Alemanha e nos EUA os efeitos são mais indiretos e complexos.

% No entanto, os modelos VAR tradicionais enfrentam limitações na identificação precisa dos choques monetários. Esta limitação motivou o desenvolvimento dos modelos SVAR, que impõem restrições teoricamente fundamentadas para identificar os choques estruturais. \citeonline{BEKAERT2013771} decompoem o indice VIX em componentes de aversão ao risco e incerteza, e utilizam um SVAR mostrando que uma política monetária expansionista reduz consistentemente a aversão ao risco após cerca de nove meses. \citeonline{li} emprega um MS-VAR (Markov-switching VAR) para examinar efeitos assimétricos do sentimento do investidor e da política monetária nos preços das ações, mostrando que os choques de sentimento têm impacto maior e mais duradouro durante períodos de baixa do mercado.

% Uma limitação importante dos modelos VAR/SVAR é sua incapacidade de incorporar grandes conjuntos de informação, o que pode levar ao problema de não-fundamentalidade dos choques estruturais. \citeonline{alessi} fornecem evidência clara deste problema ao comparar os resultados de um SVAR com um modelo de fatores dinâmicos estrutural (SDFM) usando o mesmo instrumento de alta frequência para identificar choques monetários. Os autores mostram que os SVARs subestimam consistentemente a magnitude do impacto da política monetária e produzem respostas contra-intuitivas em várias variáveis financeiras. Em contraste, o SDFM, ao incorporar um conjunto mais amplo de informações, encontra efeitos mais fortes e mais rápidos nos preços dos ativos, consistentes com a teoria macroeconômica padrão e com evidências de estudos de eventos.

% \citeonline{hanisch} utilizam um modelo dinâmico estrutural de fatores (SDFM) com 135 variáveis para analisar a eficácia da política monetária convencional e não convencional do Banco do Japão. Seus resultados mostram que choques expansionistas estimularam a economia em ambos os períodos, mas com efeitos diferentes: no período pré-QE, a redução das taxas de juros teve forte impacto na atividade econômica real, enquanto no período QE, o aumento da base monetária teve efeito mais forte nos preços dos ativos.

% Os modelos FAVAR também têm contribuído significativamente para esta literatura ao combinar a abordagem de fatores com a estrutura VAR tradicional. \citeonline{lutz} examinam como a política monetária convencional e não convencional afeta o sentimento dos investidores, encontrando que uma redução surpresa na taxa dos fed funds durante o período convencional leva a um aumento significativo no sentimento dos investidores e nos preços dos ativos. \citeonline{miranda} utilizam um FAVAR para demonstrar como um aperto monetário americano provoca uma desalavancagem significativa de intermediários financeiros globais e uma contração nos preços globais de ativos de risco.

% Metodologias mais recentes têm buscado refinar ainda mais a identificação dos choques monetários. \citeonline{miranda-transmi} desenvolvem um novo instrumento baseado em surpresas de alta frequência nos mercados futuros de juros, controlando para as projeções econômicas do Fed. Esta abordagem permite identificar com maior precisão o impacto dos choques monetários nos spreads corporativos e no prêmio de risco dos títulos. Em linha similar, \citeonline{pascal} empregam um VAR com parâmetros variando no tempo para mostrar que os preços de ações e imóveis apresentam respostas heterogêneas a choques monetários ao longo do tempo.


\section{Metodologia e Estratégia Empírica}
\label{sec:metodologia}

\subsection{O Modelo de Fatores Dinâmicos Estruturais}

O modelo empírico é um Modelo de Fatores Dinâmicos Estruturais (SDFM), representado por três equações:
\begin{align}
    \overset{n\times 1}{X_t} &= \overset{n\times r}{\Lambda} \overset{r\times 1}{F_t} + \overset{n\times 1}{e_t}\label{eq:4}\\[4pt]
    \overset{r\times r}{\Phi(L)}\overset{r\times 1}{F_t} &= \overset{r\times q}{G} \overset{q\times 1}{\eta_t} \label{eq:5}\\[4pt]
    \overset{q\times 1}{\eta_t} &= \overset{q\times q}{H} \overset{q\times 1}{\varepsilon_t} \label{eq:6}
\end{align}
A \cref{eq:4} é a equação de observação: o vetor de $n$ variáveis observáveis $X_t$ decompõe-se em um componente comum, dado por $r$ fatores latentes $F_t$ e pela matriz de cargas $\Lambda$, e um componente idiossincrático $e_t$. A \cref{eq:5} descreve a dinâmica dos fatores, com $\Phi(L) = I - \Phi_1 L - \dots - \Phi_p L^p$ um polinômio de defasagens de ordem $p$ e $\eta_t$ um vetor de $q$ inovações fatoriais. A \cref{eq:6} liga essas inovações aos $q$ choques estruturais $\varepsilon_t$ pela matriz $H$. A ordem do VAR dos fatores é fixada em $p = 6$, suficiente para acomodar a persistência das séries mensais brasileiras sem sobreparametrizar a dinâmica.

Combinando as três equações, o modelo assume a forma de média móvel que escreve as observáveis em termos dos choques estruturais:
\begin{align}\label{eq:7}
    X_t=\Lambda \Phi(L)^{-1}GH\varepsilon_t + e_t
\end{align}
A estimação requer determinar o número de fatores estáticos $r$ e dinâmicos $q$ e normalizar as matrizes $\Lambda$, $G$ e $H$.

O painel é não-estacionário por construção, pois reúne séries em nível, séries em taxa e séries já diferenciadas, com ordens de integração distintas. Extrair fatores diretamente sobre esse painel faria com que as séries de maior variância dominassem a decomposição. Seguindo \citeonline{barigozzi2016non}, cada série sem tendência $X_{it}$ é dividida pelo desvio padrão de sua primeira diferença $s_{iy}$:
\begin{equation}
    Z_{it} = \frac{X_{it}}{s_{iy}}
\end{equation}
A padronização equaliza a escala das flutuações sem diferenciar as séries, o que preserva a informação de baixa frequência usada pelo modelo. Por essa razão os critérios de \citeonline{bai-ng} não entram na forma padrão: eles pressupõem um painel estacionário, hipótese que o desenho dos dados viola. O tratamento de \citeonline{barigozzi2016non} mantém a estrutura de fatores válida sob não-estacionariedade.

O uso de um painel grande responde ao problema de não-fundamentalidade dos choques estruturais \cite{villaverde}. Um VAR de pequena dimensão raramente contém a informação que os agentes e o banco central de fato observam, e o choque recuperado nesse espaço restrito é uma combinação de choques estruturais, não o choque de política monetária isolado. Ao extrair fatores comuns de um conjunto amplo de variáveis, o SDFM aproxima o espaço de informação relevante e recupera choques fundamentais, como mostram \citeonline{alessi} ao comparar o SDFM a um VAR de referência sob o mesmo instrumento.

\subsection{Base de Dados}

A base reúne cerca de 110 séries mensais para a economia brasileira entre janeiro de 2013 e setembro de 2025, o que corresponde a 147 meses após o alinhamento das séries. O início em 2013 acompanha a mudança metodológica da PNAD Contínua no mercado de trabalho. As fontes são oficiais (Banco Central do Brasil, IBGE, B3, Tesouro Nacional e ANBIMA), somadas a indicadores de risco soberano e mercado externo (EMBI, CDS de cinco anos e MSCI) e a um índice de incerteza de política econômica (EPU). Os blocos cobrem atividade real, mercado de trabalho, inflação, commodities, agregados monetários e de crédito, e preços de ativos.

As taxas de juros de maturidade fixa (3 meses, 6 meses, 1, 2, 5 e 10 anos) vêm do ajuste do modelo de \citeonline{svensson1994estimating} sobre os contratos de DI futuro em cada data. As séries nominais entram em logaritmo e, quando os testes indicam componente sazonal, são ajustadas pelo procedimento X-13ARIMA-SEATS \cite{sax2018seasonal}. A lista completa de séries, fontes e transformações está no Anexo (Tabela \ref{tab:lista_variaveis}).

\subsection{Identificação por instrumento externo}

A identificação do choque de política monetária usa um instrumento externo, no arcabouço de proxy-SVAR de \citeonline{stockwatson2018}, \citeonline{mertensravn2013} e \citeonline{gertler2015}, transposto para o DFM como em \citeonline{alessi}. A abordagem dispensa a ordenação recursiva entre os fatores. No lugar dela entra uma variável observada fora do modelo, correlacionada com o choque de interesse e ortogonal aos demais choques estruturais.

Seja $z_t$ o instrumento e $\varepsilon_t^{p}$ o choque de política monetária da \cref{eq:6}. O instrumento é válido sob duas condições: relevância, $E[z_t \varepsilon_t^{p}] \neq 0$, e exogeneidade, $E[z_t \varepsilon_t^{j}] = 0$ para todo choque estrutural $j \neq p$ \cite{stockwatson2018}. Sob essas condições, a projeção das inovações fatoriais $\eta_t$ sobre o instrumento identifica, a menos de escala, a coluna de $H$ associada ao choque de política monetária. Empilhando $\eta_t$ e $z_t$ na amostra alinhada, essa coluna é estimada por $\hat{H}_{\cdot p} = (Z'\eta)/(Z'Z)$. O alinhamento temporal entre o instrumento mensal e as inovações do modelo segue a rotina de seleção de amostra de \citeonline{alessi}.

O choque é normalizado para elevar o yield de 6 meses em 50 pontos-base no impacto, ou $0{,}005$ em proporção decimal. A escolha do yield de 6 meses como variável de política, em vez da Selic overnight, segue a distinção de \citeonline{gertler2015} entre o instrumento de política, a taxa overnight, e o indicador de política, uma taxa de prazo mais longo que incorpora a revisão das expectativas de juros embutida na decisão. No caso brasileiro a distinção é mais aguda. As decisões do Copom saem após o fechamento do mercado, e a surpresa é medida na curva de DI entre a quarta e a quinta-feira, não na taxa overnight. A Selic mensal, uma média do overnight dentro do mês, quase não embute essa surpresa. A regressão da inovação AR(6) do yield de 6 meses sobre o instrumento primário tem F de primeiro estágio de $25{,}2$. Com a Selic mensal no lugar do yield, o F não passa de $1{,}6$ em nenhuma dimensão da varredura de especificações.

\subsection{Construção do instrumento}

O instrumento primário combina três camadas: a surpresa de juros em dias de reunião do Copom, medida como em \citeonline{gertler2015}, a ortogonalização em relação à informação pública anterior ao anúncio, de \citeonline{bauer2023}, e o filtro de sinal de \citeonline{jarocinski2020}, que separa choques de política de choques de informação.

A primeira camada mede a surpresa como a variação do DI futuro no vértice de 126 dias úteis, cerca de 6 meses, entre o fechamento da quarta-feira da reunião e o da quinta-feira seguinte. A decisão do Copom é anunciada após o fechamento do mercado, de modo que a janela de um dia captura a primeira sessão em que os preços incorporam a decisão. A amostra tem 95 reuniões entre 2013 e 2025. As surpresas são somadas dentro de cada mês, o esquema de agregação de \citeonline{jarocinski2020}, e o instrumento vale zero nos meses sem reunião.

A segunda camada remove o componente previsível da surpresa. \citeonline{bauer2023} mostram que surpresas de alta frequência são correlacionadas com informação pública anterior ao anúncio, porque o mercado subestima a resposta do banco central às notícias. Uma surpresa previsível viola a exogeneidade do instrumento. Seguindo o desenho dos autores, a surpresa é regredida sobre preditores medidos até o fechamento da quarta-feira: uma tendência temporal, as variações acumuladas em 65 pregões do Ibovespa, do S\&P 500, do VIX, do petróleo Brent, do câmbio BRL/USD e da inclinação da curva de DI (2 anos menos 3 meses), e as revisões em 20 pregões das medianas da pesquisa Focus para o IPCA acumulado em 12 meses e para a Selic no fim do ano seguinte. O resíduo dessa regressão é a surpresa ortogonalizada. O $R^2$ é $0{,}024$, na faixa esperada. \citeonline{bauer2023} obtêm entre $0{,}12$ e $0{,}20$ incluindo surpresas de divulgações macroeconômicas que não existem em frequência diária para o Brasil, e um $R^2$ alto indicaria vazamento de informação contemporânea na regressão. A exigência de que todos os preditores sejam predeterminados disciplina a especificação. Variáveis domésticas contemporâneas à janela do anúncio respondem ao próprio choque e enviesariam o resíduo. Variáveis predeterminadas, domésticas ou não, são informação disponível antes da decisão, exatamente o que a correção de \citeonline{bauer2023} usa.

A terceira camada separa o choque de política monetária do choque de informação do banco central. \citeonline{jarocinski2020} distinguem os dois pelo sinal do co-movimento entre a surpresa de juros e a do mercado de ações no mesmo dia. Juros em alta com ações em queda indicam um choque de política. Juros e ações em alta juntos indicam que o comunicado revelou informação favorável sobre a economia. O filtro classifica cada reunião pelo co-movimento dos resíduos ortogonalizados da surpresa de DI e do retorno do Ibovespa, mantém os dias de co-movimento negativo e zera os demais. Das 95 reuniões, 62 são classificadas como choque de política e 33, ou $34{,}7\%$, como informação, proporção próxima do terço que os autores encontram para o Fed. Como os preditores da segunda camada são predeterminados, a classificação preserva o co-movimento observado no dia do anúncio, em vez de depender de resíduos construídos com informação da própria janela. A série mensal resultante é o instrumento primário do trabalho.

\subsection{Seleção de $(r, q)$ e estimação}

Os critérios de \citeonline{bai-ng} na variante padronizada e o procedimento de \citeonline{amengual2007consistent} servem de referência para a dimensão do sistema e apontam $(r, q) = (5, 4)$. A escolha final é $(r, q) = (6, 5)$, definida pela força do primeiro estágio nas duas janelas amostrais. A \autoref{tab:rq_sweep} reporta, para o instrumento primário, a estatística $\xi_{mp}$ de \citeonline{montielolea} por dimensão e janela. A estatística é a Wald de primeiro estágio na direção do impacto sobre o yield de 6 meses, robusta à heterocedasticidade. Valores acima de $10$ sustentam bandas convencionais. Valores entre $3{,}84$ e $10$ indicam instrumento fraco com conjunto de confiança de Anderson-Rubin limitado.

O par $(6, 5)$ é o único acima de $10$ na janela pré-COVID, com $12{,}49$, e fica em $6{,}94$ na amostra completa, dentro da zona de instrumento fraco. A amostra completa favorece dimensões maiores, com $13{,}13$ em $(8, 8)$, mas $r \geq 7$ colapsa na janela pré-COVID, com $6{,}90$ em $(7, 6)$ e $5{,}14$ em $(8, 8)$. Os 84 meses dessa janela não sustentam a estimação de um número maior de fatores. O par $(6, 5)$ é o único com identificação utilizável nas duas janelas e por isso é mantido em ambas.

\begin{table}[htbp]
\centering
\caption{Força do instrumento primário por dimensão $(r, q)$ e janela amostral}
\label{tab:rq_sweep}
\begin{threeparttable}
\begin{tabular}{ccc}
\toprule
$(r, q)$ & Amostra completa & Pré-COVID \\
\midrule
$(5, 4)$ & 5,22  & 9,69  \\
$(6, 5)$ & 6,94  & 12,49 \\
$(7, 6)$ & 9,67  & 6,90  \\
$(8, 8)$ & 13,13 & 5,14  \\
\bottomrule
\end{tabular}
\begin{tablenotes}[flushleft]
\footnotesize
\item Estatística $\xi_{mp}$ de \citeonline{montielolea}: Wald de primeiro estágio, robusta à heterocedasticidade, na direção do impacto do choque sobre o yield de 6 meses. Acima de $10$, bandas convencionais são confiáveis. Entre $3{,}84$ e $10$, o instrumento é fraco e o conjunto de Anderson-Rubin de 95\% é um intervalo limitado. Instrumento primário, com ortogonalização pré-evento e filtro de sinal. Amostra completa: 2013--2025. Pré-COVID: 2013--2019.
\end{tablenotes}
\end{threeparttable}
\end{table}

Os fatores estáticos são extraídos por componentes principais sobre o painel padronizado. A dinâmica segue o VAR(6) da \cref{eq:5}. A estimativa pontual das funções de impulso-resposta usa os coeficientes de mínimos quadrados ordinários, fiel à implementação de \citeonline{alessi}. A incerteza vem de um \textit{bootstrap wild} \cite{goncalveskilian2004}, com 800 replicações, robusto à heterocedasticidade dos resíduos. A correção de viés de pequena amostra de \citeonline{kilian1998small} entra apenas no processo gerador do bootstrap, não na estimativa pontual. As funções de impulso-resposta cobrem horizontes de 0 a 48 meses, com bandas de 68\% e 90\%.

\subsection{Relevância do instrumento}

A validade da inferência no proxy-SVAR depende da força do instrumento \cite{montielolea}. Com instrumento fraco, o estimador de impacto é uma razão de normais correlacionadas, com viés na direção da identificação recursiva e cobertura incorreta das bandas. A alternativa nesse caso são os intervalos de Anderson-Rubin, válidos sob qualquer força do primeiro estágio.

Três estatísticas de primeiro estágio podem ser calculadas, e elas medem coisas diferentes. O F contra a inovação AR(6) do yield de 6 meses mede a relevância univariada para a leitura do choque em pontos de Selic. O F contra o resíduo do primeiro fator é a estatística usual de proxy-SVAR com uma única variável de forma reduzida. Para o instrumento primário na amostra completa, as duas valem $25{,}2$ e $2{,}8$. A diferença de quase uma ordem de grandeza não é contradição. A variação associada à política monetária carrega em fatores além do primeiro, e a estatística de um único fator não é o diagnóstico adequado para a projeção $\hat{H}_{\cdot p} = (Z'\eta)/(Z'Z)$, que usa o vetor $\eta_t$ inteiro.

O diagnóstico que governa essa projeção é a estatística $\xi_{mp}$ da \autoref{tab:rq_sweep}, a Wald de \citeonline{montielolea} na direção da normalização. Ela projeta o primeiro estágio na combinação das inovações fatoriais que determina o impacto sobre o yield de 6 meses, com erros padrão de Eicker-White e com o instrumento residualizado nas defasagens do VAR dos fatores. O conjunto de Anderson-Rubin de 95\% é um intervalo limitado se e somente se $\xi_{mp} > 3{,}84$. A implementação reproduz os números que \citeonline{montielolea} publicam para a aplicação ao choque de petróleo de Kilian, com Wald individual de $4{,}4$ e F robusta de $9{,}4$.

Na especificação de produção, $\xi_{mp}$ vale $6{,}94$ na amostra completa e $12{,}49$ na janela pré-COVID. Na amostra completa o instrumento é fraco no sentido de \citeonline{montielolea}, o conjunto de Anderson-Rubin é limitado e os resultados dessa janela são lidos com essa qualificação. Na janela pré-COVID as bandas convencionais são adequadas. A Wald conjunta sobre as $q$ inovações fatoriais é baixa nas duas janelas, com forma-F de $1{,}75$ e $2{,}74$. Wald conjunta baixa com $\xi_{mp}$ alto é o padrão esperado sob exogeneidade: a relevância se concentra na direção do choque de política, não espalhada pelo espaço dos fatores. O contraste entre as janelas repete o resultado da \autoref{tab:rq_sweep}. As observações a partir de 2020 adicionam ruído ao primeiro estágio, não informação.

\section{Resultados}

A \autoref{fig:irf_resultados} apresenta as funções de impulso-resposta para um conjunto de variáveis financeiras após o choque de política monetária contracionista. O impacto sobre o mercado de ações (IBRx-100) e sobre o mercado de dívida corporativa (IDA) é imediato e negativo, refletindo o canal do preço dos ativos. O efeito é particularmente pronunciado sobre as ações de empresas do segmento imobiliário (IMob).

No mercado cambial, a função de impulso-resposta indica uma forte sensibilidade da taxa USD/BRL, que se aprecia aproximadamente 8\% no pico da resposta. A queda no IPP, mais rápida que a do IPCA, corrobora essa dinâmica. No mercado de renda fixa, o impacto da política monetária varia ao longo da estrutura a termo, com os yields de vencimentos mais longos mostrando maior sensibilidade. A contração observada nos spreads de crédito sugere que a política monetária impacta as condições de financiamento, possivelmente alterando a percepção de risco dos agentes, em linha com \citeonline{ehrmann}.

\begin{figure}[H]
    \centering
    \caption{Funções de Impulso-Resposta a um Choque Monetário Contracionista}
    \includegraphics[width=0.8\textwidth]{"img/irf.png"}
    \label{fig:irf_resultados}
\end{figure}    


\section{Conclusão}

Este artigo quantificou os efeitos da política monetária sobre os preços dos ativos no Brasil, empregando um modelo de fatores dinâmicos estruturais (SDFM) para lidar com o problema da não-fundamentalidade. Os resultados mostram que os canais de transmissão monetária são robustos e rápidos: um choque contracionista impacta imediatamente o mercado de ações, a taxa de câmbio e a estrutura a termo dos juros, com efeitos heterogêneos entre os diferentes ativos.

As implicações para a política monetária são relevantes. A magnitude e a rapidez das respostas sugerem que o Banco Central deve considerar cuidadosamente os potenciais efeitos colaterais de suas decisões sobre a estabilidade financeira. A sensibilidade particular de setores como o imobiliário reforça a necessidade de uma análise granular dos impactos da política de juros. Como limitação e direção para pesquisa futura, a identificação dos choques poderia ser aprimorada com o uso de dados de alta frequência, como em \citeonline{alessi}, para isolar com ainda mais precisão a surpresa monetária.

% ----------------------------------------------------------
% Referências bibliográficas
% ----------------------------------------------------------
\newpage
\bibliography{./references.bib}




\begin{anexosenv}


\partanexos
\chapter{Tabela}

\begin{ThreePartTable}
\begin{TableNotes}
\item[*] \textbf{Códigos de Transformação:} 
0 = Sem transformação; 
1 = Log; 
2 = Primeira diferença; 
3 = Segunda diferença; 
4 = Dessazonalizada; 
5 = Log + Dessazonalizada; 
6 = Log + Primeira diferença; 
7 = Log + Segunda diferença; 
8 = Dessazonalizada + Primeira diferença; 
9 = Dessazonalizada + Segunda diferença; 
10 = Log + Dessazonalizada + Primeira diferença; 
11 = Log + Dessazonalizada + Segunda diferença.
\end{TableNotes}

\begin{longtable}{llll}
\caption{Lista de Variáveis Utilizadas}\label{tab:lista_variaveis} \\
\toprule
\textbf{Variável} & \textbf{Fonte} & \textbf{Grupo} & \textbf{Código} \\
\midrule
\endfirsthead

\toprule
\textbf{Variável} & \textbf{Fonte} & \textbf{Grupo} & \textbf{Código} \\
\midrule
\endhead



\bottomrule
\insertTableNotes
\endlastfoot

cambio usd & Yahoo Finance & Câmbio & 2 \\
cambio eur & Yahoo Finance & Câmbio & 2 \\
cambio gbp & Yahoo Finance & Câmbio & 2 \\
cambio ars & Yahoo Finance & Câmbio & 2 \\
cambio jpy & Yahoo Finance & Câmbio & 2 \\
cambio cny & Yahoo Finance & Câmbio & 0 \\
cambio rub & Yahoo Finance & Câmbio & 0 \\
juros selic & BCB (4189) & Monetária & 3 \\
juros cdi & BCB (4392) & Monetária & 3 \\
credito agro & BCB (22027) & Monetária & 6 \\
credito industria total & BCB (22043) & Monetária & 10 \\
credito construcao & BCB (22030) & Monetária & 11 \\
credito comercio & BCB (22036) & Monetária & 10 \\
credito transporte & BCB (22037) & Monetária & 6 \\
credito pessoa fisica & BCB (22050) & Monetária & 10 \\
spread icc juridica & BCB (27444) & Monetária & 8 \\
spread icc fisica & BCB (27445) & Monetária & 8 \\
titulo 0 25ano & Tesouro Nacional & Monetária & 2 \\
titulo 1ano & Tesouro Nacional & Monetária & 2 \\
titulo 2ano & Tesouro Nacional & Monetária & 2 \\
titulo 3ano & Tesouro Nacional & Monetária & 2 \\
titulo 5ano & Tesouro Nacional & Monetária & 2 \\
veiculos automoveis & BCB (7384) & Atividade Econômica & 5 \\
veiculos leves & BCB (7385) & Atividade Econômica & 10 \\
veiculos caminhoes & BCB (7386) & Atividade Econômica & 10 \\
veiculos onibus & BCB (7387) & Atividade Econômica & 1 \\
consumo gasolina & BCB (1393) & Atividade Econômica & 10 \\
consumo glp & BCB (1394) & Atividade Econômica & 5 \\
consumo oleo combustivel & BCB (1395) & Atividade Econômica & 6 \\
consumo oleo diesel & BCB (1396) & Atividade Econômica & 11 \\
consumo demais derivados & BCB (1397) & Atividade Econômica & 10 \\
consumo alcool & BCB (1401) & Atividade Econômica & 5 \\
consumo eletricidade comecial & BCB (1402) & Atividade Econômica & 11 \\
consumo eletricidade residencial & BCB (1403) & Atividade Econômica & 10 \\
consumo eletricidade industrial & BCB (1404) & Atividade Econômica & 10 \\
consumo eletricidade outros & BCB (1405) & Atividade Econômica & 10 \\
producao bens consumo & BCB (21865) & Atividade Econômica & 4 \\
producao bens duraveis & BCB (21866) & Atividade Econômica & 4 \\
producao bens nao duraveis & BCB (21867) & Atividade Econômica & 4 \\
producao bens capital & BCB (21863) & Atividade Econômica & 8 \\
producao transformacao & BCB (21862) & Atividade Econômica & 4 \\
capacidade instalada industria & BCB (28554) & Atividade Econômica & 2 \\
vendas varejo & BCB (1455) & Atividade Econômica & 4 \\
vendas servicos & BCB (23982) & Atividade Econômica & 8 \\
trab caged & BCB (28763) & Mercado de trabalho & 10 \\
trab pop forca trab & BCB (24378) & Mercado de trabalho & 6 \\
trab pop ocupada & BCB (28543) & Mercado de trabalho & 10 \\
trab tx desemprego & BCB (24369) & Mercado de trabalho & 8 \\
trab hrs trabalhadas industria & BCB (28556) & Mercado de trabalho & 6 \\
trab razao vagas desempregados & Calculado & Mercado de trabalho & 8 \\
price ipca & BCB (433) & Inflação & 0 \\
price igp m & BCB (189) & Inflação & 0 \\
price ipc & BCB (191) & Inflação & 0 \\
price incc & BCB (192) & Inflação & 4 \\
price inpc & BCB (188) & Inflação & 0 \\
price ipp & IPEA & Inflação & 0 \\
indice commodity agro & BCB (27575) & Commodities & 2 \\
indice commodity metal & BCB (27576) & Commodities & 2 \\
indice commodity energia & BCB (27577) & Commodities & 2 \\
trab menos de 1 mes & IBGE/SIDRA & Mercado de trabalho & 10 \\
trab de 1 mes a menos de 1 ano & IBGE/SIDRA & Mercado de trabalho & 10 \\
trab de 1 ano a menos de 2 anos & IBGE/SIDRA & Mercado de trabalho & 6 \\
trab x2 anos ou mais & IBGE/SIDRA & Mercado de trabalho & 10 \\
asset ibrx100 & B3 & Índices de Mercado & 0 \\
asset imat & B3 & Índices de Mercado & 0 \\
asset indx & B3 & Índices de Mercado & 0 \\
asset ifnc & B3 & Índices de Mercado & 0 \\
asset ifix & B3 & Índices de Mercado & 0 \\
asset imob & B3 & Índices de Mercado & 0 \\
asset MLCX & B3 & Índices de Mercado & 0 \\
asset smll & B3 & Índices de Mercado & 0 \\
retorno mensal & ANBIMA (IDA) & Índices de Mercado & 0 \\

\end{longtable}
\end{ThreePartTable}


\end{anexosenv}


\end{document}









